% =====================================================================
%  附录 F  图形 Ward 恒等式（appF_ward.tex；writer W10）
% =====================================================================
\chapter{图形 Ward 恒等式}\label{app:appF}

本附录给出第~9~章 Ward 恒等式部分所承诺的证明：单个纵向胶子的非阿贝尔 Ward 恒等式\cref{prop:method-ward-one}（\tagR，$R_{\mathrm{PT}}$）的证明，
树图层次两个纵向胶子的完整证明\cref{prop:appF-two-gluon}（\tagR，$R_{\mathrm{PT}}$），多胶子情形归纳步骤的论证纲要\cref{prop:appF-induction}（\tagA），
有限 Wilson 线的伸缩恒等式（\cref{ex:appF-telescoping}），以及带 $i0$ 的 Grammer--Yennie 分解及其幂次压低的数值表。记号同第~9~章：$S(p)=\ii(\sla p+m)/(p^2-m^2+i0)$，
夸克—胶子顶点 $\ii g\gamma^\mu t^a$，动量 $k$ 流入胶子所附着的线，$A(x)=\int\frac{\dd^4k}{(2\pi)^4}\ee^{-\ii k\cdot x}\tilde A(k)$。

% ---------------------------------------------------------------------
\section{单个纵向胶子}\label{sec:appF-one}

\paragraph{\cref{prop:method-ward-one}的证明。}
设 $\Phi$ 是一组外部物理粒子（夸克、横向极化的胶子）的场之积，经 LSZ 截腿后在壳，动量非例外。由\cref{thm:prelim-slavnov-taylor}的 $N=1$ 情形，
$\langle0|T\,\partial\cdot A^a(x)\,\Phi|0\rangle=0$。其来源是：在协变规范中 $\partial\cdot A^a=-\xi_g\,s\bar c^a$（$s$ 为 BRST 变换），真空 BRST 不变，故
$\langle T\,\partial\cdot A^a\,\Phi\rangle=\mp\xi_g\langle T\,\bar c^a\,s\Phi\rangle$；$s\Phi$ 是把 $\Phi$ 中某一个场换成其 BRST 变换之后的和：对夸克场为 $\ii g\,c^at^a\psi$，对胶子场为 $D_\mu c$。
前者以及 $D_\mu c$ 中的 $gf^{abc}A^bc^c$ 是复合算符，在外线的单粒子极点处没有极点，截腿后为零；$\partial_\mu c$ 给出终止于纵向极化 $\propto k_\mu$ 的鬼场传播子，
与物理（横向）极化缩并后为零。在动量空间中写出这一恒等式：Fourier 变换把 $\partial^\mu$ 换成 $-\ii k^\mu$，而外部胶子的完全传播子
$D^{\mu\nu}(k)=\frac{-\ii}{k^2+i0}\bigl[\frac{g^{\mu\nu}-k^\mu k^\nu/k^2}{1+\Pi(k^2)}+\xi_g\frac{k^\mu k^\nu}{k^2}\bigr]$ 的纵向部分不被重整化，
$k_\mu D^{\mu\nu}(k)=-\ii\xi_gk^\nu/(k^2+i0)$；因此 $0=-\ii k_\mu D^{\mu\nu}\mathcal M_\nu=-\xi_g\,(k\cdot\mathcal M)/(k^2+i0)$，即 $k_\mu\mathcal M^\mu_a=0$，对离壳的 $k$ 也成立。
在因子化中，这一恒等式用于\emph{近似之后}的硬部分：其余外线中进入 $H$ 的共线部分子已被换成在壳无质量的 $\hat k$ 并作了物理投影（\cref{def:method-collinear-approx}）。\qed

\paragraph{算符顶点。}若被积函数含有定义部分子分布的光线算符，则 K 胶子还可以附着在算符的 Wilson 线上；规范不变的算符的 Ward 恒等式包括这一“算符顶点项”。下面的例子说明它如何伸缩相消。

\begin{example}[有限 Wilson 线的伸缩恒等式]\label{ex:appF-telescoping}
取沿 $n_B$ 从 $0$ 到 $w=(0^+,w^-,\mathbf 0_T)=w^-n_B$ 的线段，$z(s)=sw$，$s\in[0,1]$，$\dd z\cdot A=w^-A^+\dd s$。由\cref{def:prelim-wilson}，其一阶项为
\begin{equation}\label{eq:appF-finite}
W^{(1)}[w,0]=\ii g\int_0^1\!\dd s\,w^-\,n_B\cdot\tilde A(k)\,t^a\,\ee^{-\ii s\,k\cdot w}=g\,t^a\,n_B\cdot\tilde A(k)\,\frac{1-\ee^{-\ii k^+w^-}}{k^+} ,
\end{equation}
它是 $k^+$ 的整函数（$k^+\to0$ 时趋于 $\ii gw^-t^an_B\cdot\tilde A$），不需要任何 $i0$。它等于两条半无限线之差：由 Wilson 线的规则（\cref{prop:appB-rules}），
$W^{(1)}[\infty n_B,0]=g\,t^an_B\cdot\tilde A/(k^+-i0)$，$W^{(1)}[w,\infty n_B]=-g\,t^a\ee^{-\ii k^+w^-}n_B\cdot\tilde A/(k^+-i0)$，其和
$g\,t^an_B\cdot\tilde A\,\frac{1-\ee^{-\ii k^+w^-}}{k^+-i0}$ 的分子在 $k^+=0$ 处为零，$i0$ 随之消失；用过去指向的两条线 $W[-\infty n_B,0]$ 与 $W[w,-\infty n_B]$（分母 $k^++i0$）也得到同一结果。
这是\cref{prop:method-collinear-wilson}中直 Wilson 线合并关系在一阶的显式形式。若 K 胶子的极化取 $\tilde A_\mu\to k_\mu$，则 $n_B\cdot k=k^+$，
\begin{equation}\label{eq:appF-telescope}
k^+\,\frac{1-\ee^{-\ii k^+w^-}}{k^+}\,g\,t^a=g\,t^a\bigl(1-\ee^{-\ii k\cdot w}\bigr) ,
\end{equation}
即“在 $0$ 处注入动量 $k$”减去“在 $w$ 处注入动量 $k$”——两端点的无穷小规范转动之差。对矩阵元 $\langle q(p')|\,\bar\psi(w)\gamma^+W[w,0]\psi(0)\allowbreak\,|q(p)+g(k)\rangle$ 在树图层次，
K 胶子附着在入射夸克上给出 $-g\,\ee^{\ii p'\cdot w}\,\bar u\gamma^+t^au$（由 (W2)），附着在出射夸克上给出 $+g\,\ee^{\ii(p'-k)\cdot w}\,\bar u\gamma^+t^au$，附着在 Wilson 线上由\cref{eq:appF-telescope}给出
$g\,\bigl(\ee^{\ii p'\cdot w}-\ee^{\ii(p'-k)\cdot w}\bigr)\bar u\gamma^+t^au$：三者之和为零。这就是\cref{prop:method-ward-one}中算符顶点项的作用（我们对 $m=0$ 与 $m=0.3$ 数值核对了这一恒等式）。% check: writing/work/W10/w10_appF.py
\uses{def:prelim-wilson,prop:appB-rules}
\end{example}

% ---------------------------------------------------------------------
\section{两个纵向胶子}\label{sec:appF-two}

\begin{proposition}[两个纵向胶子：树图层次的完整证明 \tagR]\label{prop:appF-two-gluon}
（$R_{\mathrm{PT}}$：在重整化微扰论中对红外正规化的部分子态严格。）（树图层次。）在 DIS 运动学中，两个共线的 K 胶子 $(k_1,\mu,a)$、$(k_2,\nu,b)$ 附着在出射的被击出夸克上，
$k_{1\mu}k_{2\nu}H^{\mu\nu}_{ab}=g^2t^bt^aH_0+\ii g^2[t^a,t^b]\,\bar u(p')\sla k_2S(p'-K)\Gamma$（$K=k_1+k_2$，精确成立）；对恰好共线的 $\hat k_i=(k_i^+,0,\mathbf 0_T)$，
\begin{equation}\label{eq:appF-two}
\begin{aligned}
\hat k_{1\mu}\hat k_{2\nu}H^{\mu\nu}_{ab}&=(n_B\cdot k_1)(n_B\cdot k_2)\,g^2\Bigl[\frac{t^bt^a}{n_B\cdot K\;n_B\cdot k_2}+\frac{t^at^b}{n_B\cdot K\;n_B\cdot k_1}\Bigr]H_0\\
&=(n_B\cdot k_1)(n_B\cdot k_2)\,W^{(2)}_{ab}[\infty n_B,0]\,H_0 :
\end{aligned}
\end{equation}
路径次序来自对易子项，这一阶没有鬼场或三胶子顶点图的贡献。对 $q\bar q\to g^*(k_1,a,\mu)\,g^*(k_2,b,\nu)$（夸克在壳，$Q=k_1+k_2$，$J_c^\sigma=\bar v(p_2)\gamma^\sigma t^cu(p_1)$），
$k_{1\mu}\mathcal A^{\mu\nu}_{ab}=\frac{g^2f^{abc}}{Q^2}(k_2^2g^{\nu\sigma}-k_2^\nu k_2^\sigma)J_{c\sigma}$；当胶子取非物理极化时鬼场项是显式的。（本命题只涉及部分子态。）
\uses{prop:method-ward-one,prop:appB-rules,def:prelim-feynman-rules,thm:prelim-slavnov-taylor}
\end{proposition}
% VF-WL-8, VF-WL-9

\begin{proof}
\emph{DIS。}被击出夸克线以动量 $p'-K$ 离开光子顶点 $\Gamma$，依次吸收两个胶子后以 $\bar u(p')$ 出射。两种次序给出
$H^{\mu\nu}_{ab}=\bar u(p')\bigl[\ii g\gamma^\mu t^aS(p'-k_1)\ii g\gamma^\nu t^b+\ii g\gamma^\nu t^bS(p'-k_2)\ii g\gamma^\mu t^a\bigr]S(p'-K)\Gamma$。
对最后吸收的胶子用 (W2)：$\bar u(p')\ii g\sla k_1S(p'-k_1)=g\bar u(p')$，于是
$k_{1\mu}k_{2\nu}H^{\mu\nu}_{ab}=g\,\bar u(p')\,\ii g\bigl(\sla k_2\,t^at^b+\sla k_1\,t^bt^a\bigr)S(p'-K)\Gamma$。把括号写成 $\sla K\,t^bt^a+\sla k_2\,[t^a,t^b]$，再对 $\sla K$ 用 (W2)
（$\bar u(p')\ii g\sla KS(p'-K)=g\bar u(p')$），即得精确式。对 $\hat k_i$：$\sla{\hat k}_2\sla{\hat K}\propto(\gamma^-)^2=0$，故
$\bar u(p')\sla{\hat k}_2(\sla p'-\sla{\hat K}+m)=\bar u(p')\sla{\hat k}_2(\sla p'+m)=2\hat k_2\cdot p'\,\bar u(p')$，而 $(p'-\hat K)^2-m^2=-2p'\cdot\hat K$；于是
$\ii\,\bar u(p')\sla{\hat k}_2S(p'-\hat K)\Gamma=\frac{\hat k_2\cdot p'}{\hat K\cdot p'}H_0=\frac{n_B\cdot k_2}{n_B\cdot K}H_0$。代入精确式：
$g^2\bigl[t^bt^a+[t^a,t^b]\frac{n_B\cdot k_2}{n_B\cdot K}\bigr]H_0=g^2\bigl[t^bt^a\frac{n_B\cdot k_1}{n_B\cdot K}+t^at^b\frac{n_B\cdot k_2}{n_B\cdot K}\bigr]H_0$，即\cref{eq:appF-two}的第一式。
未来指向线的二阶规则（\cref{prop:appB-rules}）中，靠近 $0$ 的胶子记为 $1$、较远的记为 $2$，因子为 $(\ii g)^2\frac{-\ii}{n\cdot(k_1+k_2)-i0}\frac{-\ii}{n\cdot k_2-i0}$，色矩阵按路径次序 $t^{a_2}t^{a_1}$；
对两种附着次序求和正是方括号，$-i0$ 来自被取代的末态传播子。两个胶子合并于三胶子顶点的图中，内胶子是共线的，属于 PDF 一侧而不在 $H$ 中；鬼场在这一阶不出现。
\emph{$q\bar q\to g^*g^*$。}$\mathcal A^{\mu\nu}_{ab}$ 由 $t$、$u$ 道两个夸克交换图与一个经三胶子顶点的 $s$ 道图组成。用 (W2)，$t$、$u$ 道图与 $k_{1\mu}$ 缩并后伸缩为
$\ii g^2\bar v\gamma^\nu[t^b,t^a]u=g^2f^{abc}J_c^\nu$；用 (W3)，$s$ 道图给出 $-g^2f^{abc}J_c^\nu+\frac{g^2f^{abc}}{Q^2}(k_2^2g^{\nu\sigma}-k_2^\nu k_2^\sigma)J_{c\sigma}$，
其中 $Q^\sigma Q^\nu$ 项因流守恒（$\bar v\sla Qu=0$）为零。两者相加，对易子项相消，即得所述恒等式。我们用显式振幅分别数值核对了 $t+u$ 与 $s$ 道两部分以及总和。% checks: w10_extra.py (b), w10_ward.py
若胶子 2 在壳而取非物理极化，右边不为零，正是第~9~章例子中的鬼场振幅；这与\cref{thm:prelim-slavnov-taylor}中单个纵向胶子的恒等式只在其余胶子都物理时才无鬼场项相符。
\end{proof}

这证明了\cref{prop:method-ward-tree}(b)(c)。\Cref{eq:appF-two}说明：对两个 K 胶子，阿贝尔部分（$t^bt^a$ 与 $t^at^b$ 的对称化）由程函恒等式给出 $g^2H_0$，而路径次序完全来自对易子项。

% ---------------------------------------------------------------------
\section{多胶子情形的归纳}\label{sec:appF-induction}

\begin{proposition}[多胶子情形的归纳论证 \tagA]\label{prop:appF-induction}
多个 K 胶子的 Ward 恒等式（第~9~章）的归纳步骤成立：设 $n-1$ 个 K 胶子的情形已成立，则加入第 $n$ 个 K 胶子后，附着在一切位置上之和等于 $H_0$ 乘以沿 $n_B$ 的 $n$ 阶 Wilson 线；
条件是外部部分子在壳且物理，共线动量取近似值 $\hat k_i$（而不是精确值），鬼场项相消或被吸收。（本命题只涉及部分子态。）
\uses{prop:method-ward-one,prop:appF-two-gluon}
\end{proposition}

\begin{proofsketch}
(1) 对第 $n$ 个 K 胶子用单胶子的 BRST 恒等式（\cref{prop:method-ward-one}的推导中保留其余离壳胶子）：其余 $n-1$ 个胶子离壳时，恒等式的右边出现鬼场项（Slavnov--Taylor 恒等式中带一个纵向胶子与若干任意胶子的那一式）。
(2) 第 $n$ 个胶子附着在硬部分的物理外线上的项，由 (W1)(W2) 伸缩为附着在 Wilson 线上的项；附着在其余 K 胶子线上的项由 (W3) 给出对易子，与 Wilson 线的路径次序组合
（$n=2$ 的情形即\cref{prop:appF-two-gluon}）。(3) 对恰好共线的 $\hat k_i\propto n_A$，三个 K 胶子之间的三胶子顶点为零（各项含 $n_A^2=0$ 或 $n_A\cdot\hat k=0$）。
(4) (W4) 中不是差分的部分产生鬼场项；它们必须彼此相消，或被吸收进 BRST 恰当的（A 类）算符而在物理矩阵元中消失。文献：\cite[第~11.9、11.3.2~节]{Collins2011}；
Collins--Rogers\cite{Collins:2008sg}指出朴素恒等式（只保留严格共线极限）有一个非零的领头余项。
\end{proofsketch}

\paragraph{地位与缺口。}[A]。第 (4) 步假定了所出现的非定域算符的 BRST 分类（\cite[第~11.3.2~节]{Collins2011}；\cref{rem:appJ-nonlocal}）——即具名缺口\cref{asm:appN-brst-nonlocal}（\tagC），
它在此被点名，但不列入\cref{prop:appF-induction}的“依赖”前提。附录~N 中使用的乘积形式\cref{asm:appN-ward-product}正是这一归纳若被补全所能建立的结论（其解除途径），因此也不是这里的前提。
双定域 PDF 算符的分类已由矩约化归结为定域情形（第~24~章，\tagA，先行工作见\cite{Braun:2019qtp}）。完整的证明需要多定域光线算符的 BRST 上同调。

% ---------------------------------------------------------------------
\section{Grammer--Yennie 分解及其 \texorpdfstring{$i0$}{i0}}\label{sec:appF-gy}

\paragraph{$i0$ 由被取代的传播子决定。}共线形式 $K_{\mu\nu}=k_\mu n_{B\nu}/(n_B\cdot k\mp i0)$（\cref{lem:method-gy}）中的 $i0$ 不是约定。DIS 中 K 胶子附着在出射的被击出夸克 $p'$ 上，
被取代的传播子 $(p'-k)^2+i0=-2p'^-k^++\cdots+i0$ 的 $k^+$ 极点在上半平面，故取 $-i0$（未来指向的 Wilson 线）；Drell--Yan 中附着在来自 $B$ 的入射反夸克上，
$(\bar p_B+k)^2+i0=2\bar p_B^-k^++\cdots+i0$ 的极点在下半平面，故取 $+i0$（过去指向的线；对 Drell--Yan，这一处方还须由附录~O 中单侧改写后的围道确认）。方向由被取代传播子的 $i0$ 与围道决定，
而不是由纯规范论证决定（\cref{rem:method-pure-gauge}）。软形式 $K_s$ 中的 $\pm i0$ 同理：软胶子 $k$ 附着在入射的 $A$ 共线线 $p$ 上时，被取代的传播子 $(p+k)^2-m^2+i0\simeq2p^+(k^-+i0)$ 给出 $1/(n_A\cdot k+i0)$；附着在出射线上时，$(p-k)^2-m^2+i0\simeq-2p^+(k^--i0)$ 给出 $1/(n_A\cdot k-i0)$（\cref{lem:method-gy}(b)）。对共线 PDF，两种方向给出同一个算符（\cref{prop:method-collinear-wilson}）；对 TMD 它们不同（第~14~章）。

\paragraph{历史与文献。}Grammer 与 Yennie 把 QED 中光子传播子的分子分为 G、K 两部分，K 部分正比于 $k_\mu k_\nu$，系数由相邻带电线的传播子构成，全部红外发散都在 K 光子中
（\cite{GrammerYennie1973}，如\cite{SenIndumathiChoudhury2019}（预印本）的式~(7)--(9)所转述；后来文献中的 $k\leftrightarrow-k$ 对称化系数是对原形式的修改）。软光子定理见\cite{Weinberg1965}。
带固定辅助矢量、近似动量与指定 $i0$ 的单侧形式是 Collins 的推广\cite[第~5.5.6、10.6、11.9.3~节]{Collins2011}；\cite[式~(11.4)--(11.6)]{Collins2011}须按作者的勘误表交换上下指标
（\cref{tab:appQ-errata-b}），所得 $K_{\mu\nu}=k_\mu w_{2\nu}/(k\cdot w_2-i0)$（$w_2=n_B$，$k$ 流入 $H$）正是 DIS 的未来指向处方。

\paragraph{数值上的幂次压低。}\Cref{tab:appF-gy}列出我们用显式 Dirac 矩阵算出的 G 项相对大小。% numbers: writing/work/W10/logs/w10_gy.log
共线情形取单圈 DIS 顶点的被积函数，胶子 $k=(0.6,\,0.4\lambda^2,\,0.5\lambda,\,-0.2\lambda)$
（光锥分量 $(k^+,k^-,\mathbf k_T)$）平行于靶夸克；软情形取附着在出射 $A$ 共线夸克上的软胶子与一个一般的软流 $S^\nu$。

\begin{table}[htbp]
\centering
\small
\renewcommand{\arraystretch}{1.2}
\begin{tabular}{@{}llccc@{}}
\toprule
附着 & 情形 & \multicolumn{2}{c}{$|G\text{ 项}|/|\text{全部}|$} & 拟合幂次\\
\cmidrule(lr){3-4}
& & $\lambda=10^{-2}$ & $\lambda=10^{-4}$ & \\
\midrule
共线 & 在壳无质量母夸克，横向光子 & $8.06\times10^{-6}$ & $8.06\times10^{-10}$ & $\lambda^{2.00}$\\
共线 & 在壳无质量母夸克，一般（非物理）光子极化 & $1.73\times10^{-4}$ & $1.79\times10^{-6}$ & $\lambda^{0.99}$\\
共线 & 离壳母夸克，一般旋量 & $4.73\times10^{-4}$ & $4.81\times10^{-6}$ & $\lambda^{1.00}$\\
\midrule
& & $\lambda=10^{-3}$ & $\lambda=10^{-5}$ & \\
\midrule
软 & 软区 $k\sim\lambda(1,1,1)$ & $8.48\times10^{-4}$ & $8.49\times10^{-6}$ & $\lambda^{1.00}$\\
软 & Glauber 区 $k\sim(\lambda^2,\lambda^2,\lambda)$ & $5.03\times10^{-2}$ & $5.00\times10^{-2}$ & $\lambda^{0.00}$（平台）\\
\bottomrule
\end{tabular}
\caption{Grammer--Yennie 分解中 G 项的数值大小（\cref{prop:method-gy-suppression}）。在壳无质量的母夸克配物理光子时 G 项按 $\lambda^2$ 压低（坏旋量分量为零），一般情形按 $\lambda$ 压低；
软 G 项在软区按 $\lambda$ 压低，在 Glauber 区为 $O(1)$——软近似在那里失效。}
\label{tab:appF-gy}
\end{table}
